\chapter{Radical Extensions}

\section{Solvable Groups}

\begin{definition}
    \begin{enumerate}
        \item Let $G$ be a finite group. A {\bf finite filtration} is a finite sequence
        \[
        \{e\}\leq G_{1}\leq G_{2}\leq \cdots \leq G_{n}=G
        \]
        such that $G_i\lhd G_{i+1}$ for all $i$.
        \item We say a filtration has no redundancies if $G_i\lneqq G_{i+1}$ for all $i$.
        \item We say a filtration has abelian quotients if $G_{i+1}/G_i$ is abelian for all $i$.
        \item We call $\{e\}\leq G$ the trivial filtration of $G$.
    \end{enumerate}
\end{definition}

\begin{definition}\label{def-solvgp}
    A finite group $G$ is called:
    \begin{enumerate}
        \item {\bf solvable} if $G$ has a finite filtration with abelian quotients;
        \item {\bf simple} if $G$ has no nontrivial proper normal subgroups.
    \end{enumerate}
    Thus, the only filtration with no redundancies for a simple group $G$ is
    $\{e\}\lhd G$.
\end{definition}

\begin{example}
    \begin{enumerate}
        \item $S_3$ is solvable: there is a filtration with no redundancies
        \[
        \lefteqn{\overbrace{\phantom{\{e\}\leq A_3}}^{\mathbb Z_3}}\{e\}
        \leq \underbrace{A_3\leq S_3}_{\mathbb Z_2}.
        \]

        \item $S_4$ is solvable: there is a filtration with no redundancies
        \[
        \lefteqn{\underbrace{\phantom{\{e\}\leq \langle(12)(34)\rangle}}_{\mathbb Z_2}}
        \{e\}\leq
        \overbrace{\langle(12)(34)\rangle\leq
        \lefteqn{\underbrace{\phantom{K_4\leq A_4}}_{\mathbb Z_3}}K_4}^{\mathbb Z_2}
        \leq \overbrace{A_4\leq S_4}^{\mathbb Z_2}.
        \]

        \item All finite abelian groups are solvable.

        \item $A_5$ is simple. Since $A_5$ is nonabelian, $A_5$ is not solvable.
    \end{enumerate}
\end{example}

Before entering the main content of this chapter, we highlight the idea of solving polynomial equations by radicals:
\[
    \begin{tikzcd}
        {\boxed{\text{Roots of }p(x)\in K[x]\text{ are solvable by radicals }(+,-,\times,\div,\sqrt[n]{\,})}} \\
        {\boxed{\text{The Galois group }\operatorname{Gal}(L/K)\text{ of }p(x)\text{ is solvable}}}
        \arrow[leftrightarrow, from=1-1, to=2-1]
    \end{tikzcd}
\]

\begin{proposition}\label{prop-subgpsolvable}
    Let $G$ be a solvable group.
    \begin{enumerate}
        \item If $H\leq G$, then $H$ is solvable.
        \item If $H\lhd G$, then $G/H$ is solvable.
    \end{enumerate}
\end{proposition}

\begin{proof}
    Let
    \[
    \{e\}\leq G_1\leq \cdots\leq G_n=G
    \]
    be a filtration with abelian quotients.

    \begin{enumerate}
        \item Consider
        \[
        \{e\}\leq G_1\cap H\leq \cdots\leq G_n\cap H=H. \tag{$*$}
        \]
        If $r\in G_{i+1}\cap H$ and $q\in G_i\cap H$, then
        \[
        rqr^{-1}\in G_i
        \qquad\text{and}\qquad
        rqr^{-1}\in H.
        \]
        Hence $G_i\cap H\lhd G_{i+1}\cap H$, so $(*)$ is a filtration of $H$.

        Consider the homomorphism
        \[
        \phi:G_{i+1}\cap H\longrightarrow G_{i+1}/G_i,
        \qquad
        r\longmapsto rG_i.
        \]
        Its kernel is $G_i\cap H$. Therefore, by the First Isomorphism Theorem,
        \[
        \frac{G_{i+1}\cap H}{G_i\cap H}
        \cong \operatorname{im}\phi
        \leq \frac{G_{i+1}}{G_i}.
        \]
        Since $G_{i+1}/G_i$ is abelian, so is
        $(G_{i+1}\cap H)/(G_i\cap H)$. Thus $H$ is solvable.

        \item Let
        \[
        \pi:G\longrightarrow G/H,
        \qquad
        g\longmapsto gH.
        \]
        Since $G_i\lhd G_{i+1}$, we have
        \[
        \pi(G_i)\lhd \pi(G_{i+1}).
        \]
        Hence
        \[
        \{e\}\leq \pi(G_1)\leq\cdots\leq \pi(G_n)=G/H
        \tag{$**$}
        \]
        is a filtration of $G/H$.

        The map
        \[
        G_{i+1}/G_i\longrightarrow \pi(G_{i+1})/\pi(G_i),
        \qquad
        rG_i\longmapsto \pi(r)\pi(G_i),
        \]
        is a well-defined surjective homomorphism. Therefore
        $\pi(G_{i+1})/\pi(G_i)$ is a quotient of the abelian group
        $G_{i+1}/G_i$, and hence is abelian. Thus $G/H$ is solvable.
    \end{enumerate}
\end{proof}

\begin{proposition}\label{prop-subgroupQuotientSolvableImpliesSolvable}
    Let $G$ be a finite group and $H\lhd G$. If $H$ and $G/H$ are solvable, then $G$ is solvable.
\end{proposition}

\begin{proof}
    Let
    \[
    \{e\}=H_0\leq H_1\leq\cdots\leq H_m=H
    \]
    and
    \[
    \{e\}=\overline G_0\leq \overline G_1\leq\cdots\leq \overline G_n=G/H
    \]
    be filtrations with abelian quotients of $H$ and $G/H$, respectively.

    Let $\pi:G\to G/H$ be the quotient map, and put
    \[
    G_i:=\pi^{-1}(\overline G_i).
    \]
    Then $G_0=H$, $G_n=G$, and
    \[
    H=G_0\leq G_1\leq\cdots\leq G_n=G
    \]
    is a filtration. Moreover, by the Correspondence Theorem,
    \[
    G_{i+1}/G_i\cong \overline G_{i+1}/\overline G_i,
    \]
    so these quotients are abelian. Therefore
    \[
    \{e\}=H_0\leq H_1\leq\cdots\leq H_m=H=G_0
    \leq G_1\leq\cdots\leq G_n=G
    \]
    is a filtration of $G$ with abelian quotients. Hence $G$ is solvable.
\end{proof}

\begin{example}
    \begin{enumerate}\setcounter{enumi}{4}
        \item $S_5$ is not solvable since $A_5\leq S_5$ is not solvable.
        \item For any $n\geq 5$, $A_n$ and $S_n$ are not solvable since
        $A_5\leq A_n$ and $S_5\leq S_n$.
    \end{enumerate}
\end{example}

\begin{proposition}
    Let $G$ be a group such that $|G|=p^n$ for some prime $p$. Then $G$ is solvable.
\end{proposition}

\begin{proof}
    We prove the proposition by induction on $n$.

    When $n=1$, $G\cong\mathbb Z_p$, so $G$ is solvable.

    Suppose the proposition holds for every group of order $p^i$ with $1\leq i<k$, and let $|G|=p^k$.

    \begin{claim}
        The center $Z(G)$ is nontrivial.
    \end{claim}

    \begin{proof}[Proof of claim]
        Let $G$ act on itself by conjugation:
        \[
        h\cdot g=hgh^{-1}.
        \]
        Its orbits are the conjugacy classes. By the Orbit--Stabilizer Theorem,
        every conjugacy class has cardinality a power of $p$. The conjugacy classes of size $1$ are precisely the elements of $Z(G)$.

        Write $b_i$ for the number of conjugacy classes of cardinality $p^i$. Then
        \[
        p^k=|G|=\sum_{i=0}^k b_i p^i.
        \]
        Reducing modulo $p$ gives
        \[
        b_0\equiv 0\pmod p.
        \]
        Since the identity already gives one conjugacy class of size $1$, we have $b_0\geq p$. Thus $|Z(G)|=b_0>1$.
    \end{proof}

    Since $Z(G)$ is abelian, it is solvable. Also,
    \[
    |G/Z(G)|=p^{k'}
    \]
    for some $k'<k$, so $G/Z(G)$ is solvable by induction. By Proposition
    \ref{prop-subgroupQuotientSolvableImpliesSolvable}, $G$ is solvable.
\end{proof}

\begin{definition}
    The length of a finite group $G$ is
    \[
    \operatorname{length}(G)
    =\sup\{n\mid \{e\}\leq G_1\leq\cdots\leq G_n=G
    \text{ is a filtration with no redundancies}\}.
    \]
\end{definition}

\begin{example}
    For $S_4$, there is a filtration
    \[
    \lefteqn{\underbrace{\phantom{\{e\}\leq\langle(12)(34)\rangle}}_{2}}
    \{e\}\leq
    \overbrace{\langle(12)(34)\rangle\leq
    \lefteqn{\underbrace{\phantom{K_4\leq A_4}}_{3}}K_4}^{2}
    \leq \overbrace{A_4\leq S_4}^{2},
    \]
    so $\operatorname{length}(S_4)\geq 4$.

    On the other hand, suppose
    \[
    \{e\}=G_0\lneqq G_1\lneqq\cdots\lneqq G_n=S_4
    \]
    is a filtration with no redundancies, and put $a_i=[G_i:G_{i-1}]$.
    Then every $a_i\geq 2$ and
    \[
    a_1a_2\cdots a_n=|S_4|=24=2^3\cdot 3.
    \]
    Hence $n\leq 4$. Therefore $\operatorname{length}(S_4)=4$.
\end{example}

\begin{proposition}\label{prop-maximalFiltrationCyclicGroups}
    Suppose $G$ is solvable and
    \[
    \{e\}=G_0\lneqq G_1\lneqq\cdots\lneqq
    G_{\operatorname{length}(G)}=G
    \]
    is a filtration of maximal length. Then every quotient $G_{i+1}/G_i$ is cyclic of prime order.
\end{proposition}

\begin{proof}
    Fix $i$. By maximality of the filtration, the quotient $G_{i+1}/G_i$ is simple. Indeed, if
    \[
    \{e\}\lneqq H/G_i\lneqq G_{i+1}/G_i
    \]
    were a proper nontrivial normal subgroup, then
    \[
    G_i\lneqq H\lneqq G_{i+1},
    \qquad
    G_i\lhd H\lhd G_{i+1},
    \]
    and we could insert $H$ into the filtration, contradicting maximality.

    Since $G$ is solvable, Proposition \ref{prop-subgpsolvable} implies that $G_{i+1}$ is solvable, and hence
    $G_{i+1}/G_i$ is solvable. Thus $G_{i+1}/G_i$ is both simple and solvable.

    A nontrivial finite simple solvable group must be abelian: otherwise a filtration with abelian quotients would contain a proper nontrivial normal subgroup. A finite simple abelian group is cyclic of prime order. Therefore
    \[
    G_{i+1}/G_i\cong\mathbb Z_p
    \]
    for some prime $p$.
\end{proof}

\section{Solvability by Radicals}

In this section, we assume $\operatorname{char}(K)=0$ for simplicity.

\begin{definition}\label{def-radical}
    An extension $M:K$ is called {\bf radical} if
    \begin{enumerate}
        \item $M=K(\alpha_1,\dots,\alpha_k)$;
        \item there exist $n_1,\dots,n_k\in\mathbb N$ such that
        \[
        \alpha_i^{n_i}\in K(\alpha_1,\dots,\alpha_{i-1})
        \qquad (1\leq i\leq k),
        \]
        where $K(\alpha_1,\dots,\alpha_0):=K$.
    \end{enumerate}
\end{definition}

\begin{example}
    $\mathbb Q(\sqrt{2+\sqrt[3]{2}}):\mathbb Q$ is radical since
    \[
    \mathbb Q(\sqrt{2+\sqrt[3]{2}})
    =\mathbb Q(\sqrt[3]{2},\sqrt{2+\sqrt[3]{2}}),
    \]
    and
    \[
    (\sqrt[3]{2})^3\in\mathbb Q,
    \qquad
    (\sqrt{2+\sqrt[3]{2}})^2\in\mathbb Q(\sqrt[3]{2}).
    \]
\end{example}

\begin{definition}\label{def-solvext}
    Let $p(x)\in K[x]$ be separable, and let $L:K$ be its splitting field.
    We say $L:K$ is a {\bf solvable extension} if there exists an extension
    $M:L:K$ such that $M:K$ is radical.
\end{definition}

\begin{remark}\label{rmk-radical}
    \begin{enumerate}
        \item If $L:K$ is solvable, then every root $\alpha\in L$ of $p(x)$ lies in
        $M=K(\alpha_1,\dots,\alpha_k)$, and can therefore be expressed using the operations
        $+,-,\times,\div,\sqrt[n]{\,}$ starting from elements of $K$.
        \item Cyclotomic and Kummer extensions are radical extensions.
        \item If $M':M$ and $M:K$ are radical, then $M':K$ is also radical.
    \end{enumerate}
\end{remark}

Here is our {\bf MAIN GOAL}: if $L:K$ is Galois, then
\begin{center}
    \fbox{$\operatorname{Gal}(L/K)\text{ is solvable}
    \iff L:K\text{ is a solvable extension.}$}
\end{center}

We begin with the $\Rightarrow$ direction.

\begin{lemma}\label{lem-galclosure}
    Let $P:K$ be a finite extension. Then there exists an extension $J:P$ such that $J:K$ is Galois.
\end{lemma}

\begin{proof}
    Since $\operatorname{char}(K)=0$, the Primitive Element Theorem gives
    \[
    P=K(\delta)
    \]
    for some $\delta\in P$. Let $m_\delta(x)\in K[x]$ be the minimal polynomial of $\delta$, and let $J$ be its splitting field in an algebraic closure containing $P$.
    Then $J:K$ is Galois and $P=K(\delta)\subseteq J$.
\end{proof}

\begin{theorem}\label{thm-radleft}
    Let $L:K$ be Galois with
    \[
    |\operatorname{Gal}(L/K)|=m.
    \]
    Suppose $\operatorname{Gal}(L/K)$ is solvable. Let $\omega$ be a primitive $m$-th root of unity. Then there exists an extension $P:L$ such that:
    \begin{enumerate}
        \item $P:K$ is Galois;
        \item $P$ contains the cyclotomic extension $K(\omega)$;
        \item $P:K(\omega)$ is a composition of Kummer extensions.
    \end{enumerate}
\end{theorem}

\begin{proof}
    Set
    \[
    P:=L(\omega).
    \]
    Since $L:K$ and $K(\omega):K$ are Galois, their compositum $P:K$ is Galois. This proves (1), while (2) is immediate from the definition of $P$.

    Consider the restriction map
    \[
    \rho:\operatorname{Gal}(P/K(\omega))
    \longrightarrow \operatorname{Gal}(L/K),
    \qquad
    \sigma\longmapsto \sigma|_L.
    \]
    It is well-defined because $L:K$ is Galois, so every $K$-automorphism of $P$ preserves $L$. It is injective: if $\sigma|_L=\operatorname{id}_L$, then $\sigma$ fixes both $L$ and $\omega$, hence fixes $P=L(\omega)$.

    Thus $\operatorname{Gal}(P/K(\omega))$ is a subgroup of the solvable group $\operatorname{Gal}(L/K)$, and is therefore solvable by Proposition \ref{prop-subgpsolvable}. By Proposition \ref{prop-maximalFiltrationCyclicGroups}, there is a filtration
    \[
    \{e\}=H_0\lhd H_1\lhd\cdots\lhd H_r
    =\operatorname{Gal}(P/K(\omega))
    \]
    such that
    \[
    H_{i+1}/H_i\cong\mathbb Z_{q_i}
    \]
    for some prime $q_i$.

    Let
    \[
    P_i:=P^{H_i}.
    \]
    By the Fundamental Theorem of Galois Theory,
    \[
    P=P_0\supseteq P_1\supseteq\cdots\supseteq P_r=K(\omega),
    \]
    and, since $H_i\lhd H_{i+1}$,
    \[
    \operatorname{Gal}(P_i/P_{i+1})
    \cong H_{i+1}/H_i
    \cong\mathbb Z_{q_i}.
    \]

    Since every $q_i$ divides
    \[
    |\operatorname{Gal}(P/K(\omega))|,
    \]
    and this group injects into $\operatorname{Gal}(L/K)$, we have $q_i\mid m$. Therefore $K(\omega)$, and hence every $P_{i+1}$, contains all $q_i$-th roots of unity. Since
    \[
    \operatorname{Gal}(P_i/P_{i+1})\cong\mathbb Z_{q_i},
    \]
    Theorem \ref{thm-cyclicGalImpliesKummerExtension} implies that each
    \[
    P_i:P_{i+1}
    \]
    is a Kummer extension. Hence $P:K(\omega)$ is a composition of Kummer extensions.
\end{proof}

\begin{corollary}
    Let $L:K$ be Galois. If $\operatorname{Gal}(L/K)$ is solvable, then $L:K$ is a solvable extension.
\end{corollary}

\begin{proof}
    By Theorem \ref{thm-radleft}, there is an extension $P:L$ such that
    $P:K(\omega)$ is a composition of Kummer extensions and $K(\omega):K$ is cyclotomic. By Remark \ref{rmk-radical}(2), both cyclotomic and Kummer extensions are radical, and by Remark \ref{rmk-radical}(3), a composition of radical extensions is radical. Hence $P:K$ is radical, so $L:K$ is solvable.
\end{proof}

We now prove the $\Leftarrow$ direction. We first need a slightly stronger version of Lemma \ref{lem-galclosure}.

\begin{lemma}\label{lem-radgalois}
    Let $P:K$ be a radical extension. Then there exists an extension $P':P$ such that $P':K$ is both radical and Galois.
\end{lemma}

\begin{proof}
    Write
    \[
    P=K(\alpha_1,\dots,\alpha_k),
    \qquad
    \alpha_i^{n_i}\in K(\alpha_1,\dots,\alpha_{i-1}).
    \]
    Let $J:P:K$ be a Galois closure as in Lemma \ref{lem-galclosure}, and write
    \[
    \operatorname{Gal}(J/K)=\{\sigma_1=e,\sigma_2,\dots,\sigma_r\}.
    \]
    Define
    \[
    P':=K\bigl(\sigma_j(\alpha_i)\mid
    1\leq j\leq r,\ 1\leq i\leq k\bigr).
    \]
    Since $\sigma_1=e$, we have $P\subseteq P'$.

    We claim first that $P':K$ is radical. Order the generators block by block:
    \[
    \sigma_1(\alpha_1),\dots,\sigma_1(\alpha_k),
    \sigma_2(\alpha_1),\dots,\sigma_2(\alpha_k),\dots.
    \]
    For every $j$ and $i$,
    \[
    \bigl(\sigma_j(\alpha_i)\bigr)^{n_i}
    =\sigma_j(\alpha_i^{n_i})
    \in K\bigl(\sigma_j(\alpha_1),\dots,
    \sigma_j(\alpha_{i-1})\bigr),
    \]
    which is contained in the field generated by the preceding elements in this ordering. Thus $P':K$ is radical.

    Finally, for every $\tau\in\operatorname{Gal}(J/K)$, left multiplication by $\tau$ permutes the elements $\sigma_j$ of $\operatorname{Gal}(J/K)$. Hence
    \[
    \tau(P')=P'.
    \]
    By Remark \ref{rmk-GaloisExtensionEquiFixedField}, $P':K$ is Galois.
\end{proof}

\begin{theorem}
    Let $L:K$ be Galois. If $L:K$ is a solvable extension, then
    $\operatorname{Gal}(L/K)$ is a solvable group.
\end{theorem}

\begin{proof}
    By hypothesis, there is an extension $P:L$ such that $P:K$ is radical. By Lemma \ref{lem-radgalois}, there exists $P':P$ such that $P':K$ is both radical and Galois.

    Write
    \[
    P'=K(\alpha_1,\dots,\alpha_k),
    \qquad
    \alpha_i^{n_i}\in K(\alpha_1,\dots,\alpha_{i-1})
    \quad (1\leq i\leq k).
    \]
    Put
    \[
    t:=\prod_{i=1}^k n_i,
    \]
    let $\omega$ be a primitive $t$-th root of unity, and set
    \[
    E:=P'(\omega).
    \]
    Since $P':K$ and $K(\omega):K$ are Galois, their compositum $E:K$ is Galois.

    For $0\leq i\leq k$, define
    \[
    E_i:=K(\omega)(\alpha_1,\dots,\alpha_i),
    \qquad
    E_0=K(\omega).
    \]
    Then
    \[
    E_0\subseteq E_1\subseteq\cdots\subseteq E_k=E.
    \]
    Since $n_i\mid t$, the polynomial $x^{n_i}-1$ splits in $E_{i-1}$. Moreover,
    \[
    \alpha_i^{n_i}\in K(\alpha_1,\dots,\alpha_{i-1})
    \subseteq E_{i-1},
    \]
    so $\alpha_i$ is a root of
    \[
    x^{n_i}-\alpha_i^{n_i}\in E_{i-1}[x].
    \]
    Because $E_{i-1}$ contains all $n_i$-th roots of unity, adjoining $\alpha_i$
    adjoins every root of $x^{n_i}-\alpha_i^{n_i}$. Hence $E_i:E_{i-1}$
    is Galois. By Lemma \ref{lem-kummer},
    \[
    \operatorname{Gal}(E_i/E_{i-1})\hookrightarrow \mu_{n_i}(E_{i-1}),
    \]
    so $\operatorname{Gal}(E_i/E_{i-1})$ is cyclic, and in particular abelian.
    Thus $E_i:E_{i-1}$ is a Kummer extension (of degree dividing $n_i$).
    Put
    \[
    G_i:=\operatorname{Gal}(E/E_i),
    \qquad 0\leq i\leq k.
    \]
    Since $E_i:E_{i-1}$ is Galois, the Fundamental Theorem of Galois Theory gives
    \[
    G_i\lhd G_{i-1}
    \qquad\text{and}\qquad
    G_{i-1}/G_i\cong\operatorname{Gal}(E_i/E_{i-1}),
    \]
    so every quotient $G_{i-1}/G_i$ is abelian. Also,
    \[
    G_k=\{e\},
    \]
    and, because $K(\omega):K$ is cyclotomic,
    \[
    \frac{\operatorname{Gal}(E/K)}{G_0}
    \cong \operatorname{Gal}(K(\omega)/K)
    \]
    is abelian. Thus
    \[
    \{e\}=G_k\lhd G_{k-1}\lhd\cdots\lhd G_0
    \lhd \operatorname{Gal}(E/K)
    \]
    is a filtration with abelian quotients. Hence $\operatorname{Gal}(E/K)$ is solvable.

    Finally, $L\subseteq E$ and $L:K$ is Galois, so
    \[
    \operatorname{Gal}(L/K)
    \cong
    \operatorname{Gal}(E/K)/\operatorname{Gal}(E/L).
    \]
    Therefore $\operatorname{Gal}(L/K)$ is solvable by Proposition \ref{prop-subgpsolvable}(2).
\end{proof}

\section{Cubic Polynomials}

We now give a general solution to any irreducible cubic equation
\[
x^3+ax^2+bx+c=0.
\]
After the change of variable $x\mapsto x-a/3$, it is enough to consider a depressed cubic
\[
f(x)=x^3+px+q\in\mathbb Q[x].
\]
Let $L:\mathbb Q$ be its splitting field. Since $f$ is irreducible, transitivity implies
\[
\operatorname{Gal}(L/\mathbb Q)=A_3
\qquad\text{or}\qquad
S_3.
\]

Let $\alpha_1,\alpha_2,\alpha_3\in L$ be the three roots of $f(x)$, so
\[
\alpha_1+\alpha_2+\alpha_3=0,
\]
and let $\omega$ be a primitive cube root of unity. Put
\[
P=L(\omega),
\]
and define
\[
\beta:=\alpha_1+\omega\alpha_2+\omega^2\alpha_3,
\qquad
\gamma:=\alpha_1+\omega^2\alpha_2+\omega\alpha_3.
\]
Since both $L:\mathbb Q$ and $\mathbb Q(\omega):\mathbb Q$ are Galois, their compositum $P:\mathbb Q$ is Galois.

We now compute $\beta^3+\gamma^3$ and $\beta^3\gamma^3$. First,
\[
\beta+\gamma
=2\alpha_1+(\omega+\omega^2)(\alpha_2+\alpha_3)
=3\alpha_1,
\]
since $\omega+\omega^2=-1$ and $\alpha_2+\alpha_3=-\alpha_1$. Also,
\begin{align*}
\beta\gamma
&=\alpha_1^2+\alpha_2^2+\alpha_3^2
-(\alpha_1\alpha_2+\alpha_1\alpha_3+\alpha_2\alpha_3)\\
&=-3p,
\end{align*}
where we used
\[
\alpha_1\alpha_2+\alpha_1\alpha_3+\alpha_2\alpha_3=p
\]
and
\[
\alpha_1^2+\alpha_2^2+\alpha_3^2=-2p.
\]
Therefore
\begin{align*}
\beta^3+\gamma^3
&=(\beta+\gamma)^3-3\beta\gamma(\beta+\gamma)\\
&=27(\alpha_1^3+p\alpha_1)\\
&=-27q,
\end{align*}
and
\[
\beta^3\gamma^3=(\beta\gamma)^3=-27p^3.
\]
Thus $\beta^3$ and $\gamma^3$ are the roots of
\[
T^2+27qT-27p^3=0.
\]
Hence
\[
\boxed{
\beta^3,\gamma^3
=\frac{-27q\pm\sqrt{729q^2+108p^3}}{2}.}
\]
Choose cube roots $\beta,\gamma$ of these two numbers compatibly so that
\[
\beta\gamma=-3p.
\]
Then
\[
\alpha_1=\frac{\beta+\gamma}{3},
\qquad
\alpha_2=\frac{\beta\omega^2+\gamma\omega}{3},
\qquad
\alpha_3=\frac{\beta\omega+\gamma\omega^2}{3}.
\]
This gives the roots of the cubic by radicals.

\section{Insolvability of Quintic Polynomials}

Let $f(x)\in K[x]$ be irreducible of degree $5$, and let $L:K$ be its splitting field. Then
\[
\operatorname{Gal}(L/K)\leq S_5.
\]
To show that $f(x)$ is not solvable by radicals, it is enough to find an example for which the Galois group is not solvable; in particular, it is enough to find an example with Galois group $S_5$.

\begin{theorem}
Let $f(x)\in\mathbb Q[x]$ be an irreducible polynomial of prime degree $p$. If $f(x)$ has exactly $p-2$ real roots, then the splitting field $L$ of $f(x)$ satisfies
\[
\operatorname{Gal}(L/\mathbb Q)\cong S_p.
\]
\end{theorem}

\begin{proof}
    Let
    \[
    G:=\operatorname{Gal}(L/\mathbb Q)\leq S_p.
    \]
    Since $f(x)$ is irreducible of degree $p$, for any root $\alpha$ of $f$ we have
    \[
    [\mathbb Q(\alpha):\mathbb Q]=p.
    \]
    As
    \[
    \mathbb Q\subseteq\mathbb Q(\alpha)\subseteq L,
    \]
    the Tower Law gives
    \[
    p\mid [L:\mathbb Q]=|G|.
    \]
    By Cauchy's Theorem, $G$ contains an element $\sigma$ of order $p$. Since $G\leq S_p$, such an element must be a $p$-cycle.

    By hypothesis, $f(x)$ has exactly two nonreal roots. Since $f$ has real coefficients, these roots are complex conjugates. Complex conjugation preserves the splitting field $L$, fixes $\mathbb Q$, fixes the $p-2$ real roots, and interchanges the two nonreal roots. Hence its restriction to $L$ gives a transposition
    \[
    \tau\in G.
    \]

    It remains to show that a $p$-cycle together with a transposition generates $S_p$. Relabel the roots so that
    \[
    \sigma=(0\ 1\ \cdots\ p-1),
    \qquad
    \tau=(0\ d)
    \]
    for some $d\not\equiv 0\pmod p$. Then
    \[
    \sigma^k\tau\sigma^{-k}=(k\ k+d)
    \qquad (k\in\mathbb Z/p\mathbb Z).
    \]
    Since $p$ is prime, $d$ generates the additive group $\mathbb Z/p\mathbb Z$. Thus the graph with edges
    \[
    \{k,k+d\}
    \]
    is connected. The transpositions corresponding to the edges of a connected graph generate the full symmetric group. Hence
    \[
    \langle\sigma,\tau\rangle=S_p.
    \]
    Since $\sigma,\tau\in G\leq S_p$, we conclude that $G=S_p$.
\end{proof}

\begin{example}
Consider
\[
f(x)=x^5-6x+3.
\]
By Eisenstein's criterion at $3$, $f(x)$ is irreducible over $\mathbb Q$.

We claim that $f$ has exactly three real roots. First,
\[
f(-2)<0<f(-1),
\qquad
f(0)>0>f(1),
\qquad
f(1)<0<f(2),
\]
so the Intermediate Value Theorem gives at least three real roots. On the other hand,
\[
f'(x)=5x^4-6
\]
has exactly two real zeros. By Rolle's Theorem, $f$ can therefore have at most three real roots. Hence $f$ has exactly three real roots.

The theorem therefore gives
\[
\operatorname{Gal}(L/\mathbb Q)\cong S_5.
\]
Since $S_5$ is not solvable, $f(x)$ is not solvable by radicals.
\end{example}
